\documentclass{article}

% ============================================================
% MA 213 In-Class Worksheet Template
% Student-facing worksheet for lecture activities.
% ============================================================

\usepackage[margin=1in]{geometry}
\usepackage{amsmath}
\usepackage{xcolor}
\usepackage{parskip}
\usepackage{array}
\usepackage{enumitem}
\usepackage{booktabs}

\newcommand{\coursenumber}{MA 213}
\newcommand{\weekplaceholder}{13}
\newcommand{\lectureplaceholder}{32}
\newcommand{\lecturetitle}{Finding the MLE by Grid Search: Think/Pair/Share}

\newcommand{\nameblank}{\rule{1.75in}{0.4pt}}
\newcommand{\idblank}{\rule{1.55in}{0.4pt}}
\newcommand{\shortblank}{\rule{1in}{0.4pt}}
\newcommand{\longblank}{\rule{0.93\linewidth}{0.4pt}}

\begin{document}

\begin{center}
  {\LARGE \coursenumber{} Week \weekplaceholder{}: Lecture \lectureplaceholder{}}\\
  {\large \lecturetitle{}}
\end{center}

\begin{enumerate}[label=\arabic*.,leftmargin=*,itemsep=.7em]
  \item \textbf{Your name:} \nameblank \hfill \textbf{BU ID:} \idblank
\end{enumerate}

\textbf{Big idea:} You just watched the likelihood get evaluated at a grid of candidate $\mu$ values for the temperature example. Now you'll do the same thing yourself, for a \textbf{Geometric} distribution (Module 2).

\section*{The Setup}

A basketball player takes free throws until she makes her first basket. It takes her $X=4$ attempts. Recall the Geometric pmf: $P(X=x) = (1-p)^{x-1}p$. So the likelihood here is

\[
L(p) = (1-p)^{3}\,p.
\]

\section*{Step 1: Think (4 mins)}

\textbf{(a)} Fill in the table below:

\begin{center}
\renewcommand{\arraystretch}{1.5}
\begin{tabular}{c c}
\toprule
$p$ & $(1-p)^3 p$ \\
\midrule
0.10 & \\
0.20 & \\
0.25 & \\
0.30 & \\
0.40 & \\
0.50 & \\
\bottomrule
\end{tabular}
\end{center}

\textbf{(b)} Which value of $p$ in your table \textbf{maximizes} the likelihood? \shortblank

\textbf{(c)} Based on your answer, guess a \textbf{general formula} for $\hat{p}$ in terms of $X$ (the number of attempts until the first basket).
\vspace{0.3in}

\section*{Step 2: Pair (4 mins)}

\begin{enumerate}[label=\arabic*.,leftmargin=*,itemsep=.7em, start=2]
  \item \textbf{Partner's name:} \nameblank \hfill \textbf{BU ID:} \idblank
\end{enumerate}

Compare your table and your answers to (b) and (c) with your partner. Then discuss:

\begin{itemize}
  \item Why does it make intuitive sense that a \emph{larger} number of attempts before the first basket ($X$) leads to a \emph{smaller} estimated $p$?
  \item \textbf{Bonus:} using your formula from (c), predict the MLE if she'd instead made her first basket on attempt $X=6$. Check your prediction by computing $L(p)$ at $p=1/6$ and one neighboring value (e.g.\ $p=0.2$).
\end{itemize}

\textbf{Our thoughts:}
\vfill

\end{document}
